\section{Primitive boundaries of cell fans}
\label{sec:area_law_fan_frontiers}

The local mesh argument of \cite[Section~11]{OpenAI2026AreaLaw} tracks the
endpoints of boundary segments. Fix an origin, a dyadic exponent, an
integer cell index, and optional midpoint subdivisions. Write $Q$ for the
half-open square, $c$ for its center, $\mathcal E$ for its cyclic fan slots,
$P_e$ for the corresponding closed triangle, and $v_e$ for its final base
endpoint.

% Source: 10-geometry.tex, prop:two-families, lines 299--323,
% and geometry:initial-stars, lines 333--370, especially 352--359;
% revision adc7f1241b42e322a6451854ab7e4b4c146bf78a.
% This is a primitive frontier inclusion, not the active-star assertion.

% Source: 10-geometry.tex, prop:two-families, lines 308--323 and 545--559;
% revision adc7f1241b42e322a6451854ab7e4b4c146bf78a.

\begin{lemma}[Boundaries in a finite closed cover]
    \label{thm:area_law_finite_cover_frontier}
    \lean{TNLean.PEPS.AreaLaw.Geometry.finiteCover_frontier_subset}
    \leanok
    Let $Q\subseteq\mathbb R^2$, let $\mathcal E$ be a finite index set,
    and let $P_e$ be closed sets satisfying
    $\bigcup_{e\in\mathcal E}P_e=\overline Q$.
    In the finite-cover inclusion entering the local geometric argument
    of \cite[Section~11]{OpenAI2026AreaLaw}, every $e\in\mathcal E$ satisfies
    \begin{align}
        \partial P_e
         & \subseteq\partial Q\cup
        \bigcup_{\substack{f\in\mathcal E \\f\ne e}}(P_e\cap P_f).
    \end{align}
\end{lemma}
\begin{proof}\leanok
    Let $x\in\partial P_e$ with $x\notin\partial Q$.
    Since $P_e$ is closed, $x\in P_e\subseteq\overline Q$, so
    $x\in\operatorname{int}(Q)$.
    If $x$ belonged to no $P_f$ with $f\ne e$, the complement of
    $\bigcup_{f\ne e}P_f$ would be an open neighborhood of $x$ whose
    intersection with $\operatorname{int}(Q)$ lies in $P_e$.
    That would place $x$ in $\operatorname{int}(P_e)$, contradicting
    $x\in\partial P_e$.
\end{proof}

\begin{lemma}[Fan boundaries lie on outer or radial edges]
    \label{thm:area_law_fan_frontier_primitives}
    \lean{TNLean.PEPS.AreaLaw.Geometry.cellFanPolygon_frontier_subset_outer_and_radials}
    \leanok
    \uses{def:area_law_cell_fan}
    For every fan slot $e\in\mathcal E$,
    \begin{align}
        \partial P_e & \subseteq\partial Q\cup
        \bigcup_{f\in\mathcal E}[c,v_f].
    \end{align}
\end{lemma}
\begin{proof}\leanok
    \uses{thm:area_law_finite_cover_frontier,
        thm:area_law_cell_fan_cover,
        thm:area_law_fan_triangle_contact,thm:area_law_cell_fan_intersections}
    Because the closed fan triangles $P_e$ have union $\overline Q$,
    Lemma~\ref{thm:area_law_finite_cover_frontier} places every boundary
    point of $P_e$ outside $\partial Q$ in some intersection $P_e\cap P_f$
    ($f\ne e$). Any point in $P_e\cap P_f$ is either the center $c$ or a
    second common point distinct from $c$; in the latter case the
    triangle-contact classification places it on a radial segment $[c,v_f]$,
    while $c$ itself belongs to every such segment.
\end{proof}

The inclusion retains all radial edges of the fan, including seams that
disappear when neighboring triangles of the same region are merged.

